% =====================================================================
% Seção de metodologia e validação — artigo sparse-joint networks
%
% Mantida no repositório, ao lado do código que produz os números:
%   validation.py  gera results/validation/*.json
%   report.py      gera results/validation/paper/{tables,figuras}
%
% Para usar no artigo, copie as tabelas e figuras geradas para figs/ e
% faça % =====================================================================
% Seção de metodologia e validação — artigo sparse-joint networks
%
% Mantida no repositório, ao lado do código que produz os números:
%   validation.py  gera results/validation/*.json
%   report.py      gera results/validation/paper/{tables,figuras}
%
% Para usar no artigo, copie as tabelas e figuras geradas para figs/ e
% faça % =====================================================================
% Seção de metodologia e validação — artigo sparse-joint networks
%
% Mantida no repositório, ao lado do código que produz os números:
%   validation.py  gera results/validation/*.json
%   report.py      gera results/validation/paper/{tables,figuras}
%
% Para usar no artigo, copie as tabelas e figuras geradas para figs/ e
% faça % =====================================================================
% Seção de metodologia e validação — artigo sparse-joint networks
%
% Mantida no repositório, ao lado do código que produz os números:
%   validation.py  gera results/validation/*.json
%   report.py      gera results/validation/paper/{tables,figuras}
%
% Para usar no artigo, copie as tabelas e figuras geradas para figs/ e
% faça \input{methodology} depois da seção de trabalhos relacionados.
%
% Requer no preâmbulo (já presentes no template): graphicx, booktabs,
% tikz (arrows.meta, positioning), amsmath.
% =====================================================================

\section{The Sparse-Joint Core}
\label{sec:architecture}

A sparse-joint network replaces the dense bottleneck of a conventional encoder
with a \emph{graded program}: a representation that is simultaneously discrete,
as a tuple of per-group symbols, and continuous, as a hypervector that the
network can differentiate through. This section states the construction
precisely; Section~\ref{sec:protocol} describes how we test it.

\subsection{Groups and the scaffold}

Let $U$ be the number of hidden units per layer, $L$ the number of layers and
$G$ the number of \emph{groups}. A group is a sparse population of units
distributed across layers, fixed before training by a binary mask
$M \in \{0,1\}^{L \times G \times U}$, where $M_{l,g,u}=1$ means that group $g$
owns unit $u$ in layer $l$. We call $M$ the \emph{scaffold}.

Two groups are adjacent when they physically share a unit:
\begin{equation}
  (i,j) \in E \iff \exists\, l,u : M_{l,i,u} = M_{l,j,u} = 1 .
  \label{eq:edges}
\end{equation}
The edge set $E$ is therefore not a hyperparameter chosen alongside the
architecture; it is a \emph{consequence} of the scaffold, readable from the
weights themselves. This is what later allows $E$ to act as a compatibility
certificate between two independently trained networks: two networks with
isomorphic scaffolds expose the same relational skeleton, even though neither
has seen the other's data. Figure~\ref{fig:glue}(a) shows the graph used
throughout this work, for $G=8$.

Unit-to-unit transitions are routed through the scaffold: a connection between
units $u$ and $u'$ exists only if some group contains both. A group therefore
cannot read an arbitrary part of the hidden state, which is what makes the
sparsity structural rather than merely a regulariser.

\subsection{From activations to a program}

Write $h^{(l)} \in \mathbb{R}^{U \times W}$ for the hidden state at layer $l$,
each unit carrying a width-$W$ vector, with
$h^{(0)} = \mathrm{GELU}(W_{\text{in}}\, \phi(x))$ for an encoder $\phi$. Each
group pools the units it owns, across all layers:
\begin{equation}
  f_g = \frac{1}{L}\sum_{l=1}^{L}
        \frac{\sum_{u} M_{l,g,u}\, h^{(l)}_u}{\sum_{u} M_{l,g,u}},
  \qquad
  v_g = \frac{W_{\text{val}} f_g}{\lVert W_{\text{val}} f_g \rVert}.
  \label{eq:pool}
\end{equation}

Each group owns a fixed bipolar \emph{role} vector $r_g \in \{-1,+1\}^{D}$,
drawn once at initialisation. Binding is the Hadamard product, the standard
multiply-add vector-symbolic algebra: $r_g \odot v_g$ binds the group's value to
its identity, and because $r_g \odot r_g = \mathbf{1}$, binding again by the
same role unbinds it. Unary messages are written back to every unit the group
owns, and pairwise messages $r_i v_i \odot r_j v_j$ only to the units that $i$
and $j$ physically share, so interaction is confined to $E$ by construction.
The stack is then re-pooled, giving refined values $v'_g$.

The refined values are superposed into a single hypervector, and each group is
recovered from the superposition by unbinding with its own role:
\begin{equation}
  b = \frac{1}{\sqrt{G}}\sum_{g=1}^{G} r_g \odot v'_g,
  \qquad
  \tilde{v}_g = \sqrt{G}\, \bigl( b \odot r_g \bigr).
  \label{eq:bundle}
\end{equation}
Recovery is approximate: the other $G-1$ terms survive as noise whose magnitude
falls with $D$. We measure this directly (Section~\ref{sec:protocol},
\emph{VSA recovery}) rather than assuming it, because every downstream claim
about reading codes back out depends on it.

\subsection{Discretisation}

Each group owns a codebook $C_g \in \mathbb{R}^{K_g \times D}$ with its own
alphabet size $K_g$, so groups need not share a vocabulary. Cleanup is a cosine
lookup against that codebook,
\begin{equation}
  \ell_{g,c} = \frac{1}{\tau}
  \left\langle \frac{\tilde{v}_g}{\lVert \tilde{v}_g \rVert},
               \frac{C_{g,c}}{\lVert C_{g,c} \rVert} \right\rangle
  + m_{g,c},
  \label{eq:cleanup}
\end{equation}
where $m_{g,c} = -\infty$ for $c \geq K_g$ masks out fillers outside the group's
alphabet. The discrete code is $z_g = \arg\max_c \ell_{g,c}$, and gradients pass
through a straight-through estimator~\cite{bengio2013estimating,oord2017neural}, so the forward pass is exactly discrete
while the backward pass stays differentiable.

The exported \emph{program} is graded, concatenating a unary and a pairwise
bundle:
\begin{equation}
  \begin{aligned}
    P(x) &= \bigl[\, p_{\text{unary}} \;;\; p_{\text{pair}} \,\bigr],\\
    p_{\text{unary}} &= \tfrac{1}{\sqrt{G}}\textstyle\sum\nolimits_g r_g \odot \hat{v}_g,\\
    p_{\text{pair}} &= \tfrac{1}{\sqrt{|E|}}\textstyle\sum\nolimits_{(i,j) \in E}
       \hat{h}_i \odot \hat{h}_j,
  \end{aligned}
  \label{eq:program}
\end{equation}
where $\hat{h}_g = r_g \odot \hat{v}_g$ is the bound contribution of group $g$,
and $\hat{v}_g$ the quantised value of group $g$. The first term says what each
part is; the second says which parts co-occur, and only along scaffold edges.
Figure~\ref{fig:method} traces the whole path.

\subsection{Redundancy between groups}
\label{sec:penalty}

Nothing in Equations~\ref{eq:pool} to~\ref{eq:program} forces groups to encode
\emph{different} things. Trained only to classify, every group independently
converges on the label: the codes become near-copies of one another, and the
pairwise term of Equation~\ref{eq:program} carries nothing the unary term does
not already carry. We confirm this empirically in Section~\ref{sec:results}.

We therefore add a penalty on the mutual information between the soft code
distributions of every pair of groups. Writing $p_g \in \Delta^{K_g}$ for the
posterior of group $g$ on a batch, and $\hat{p}_{ij}$ for the empirical joint,
\begin{equation}
  \mathcal{L}_{\text{red}} =
  \frac{1}{G(G-1)} \sum_{i \neq j}
  \sum_{c,c'} \hat{p}_{ij}(c,c')
  \log \frac{\hat{p}_{ij}(c,c')}{\hat{p}_i(c)\,\hat{p}_j(c')},
  \label{eq:redundancy}
\end{equation}
added to the task loss with weight $\lambda$. Minimising it pushes the groups
towards statistically independent codes. The full objective is
$\mathcal{L} = \mathcal{L}_{\text{task}} + \mathcal{L}_{\text{VQ}}
+ \mathcal{L}_{\text{support}} + \lambda \mathcal{L}_{\text{red}}$.

% =====================================================================

\section{Experimental Protocol}
\label{sec:protocol}

The claims this architecture makes are about generality and about structure,
not only about accuracy, so the protocol has to be comparable across domains
and has to measure structure explicitly.

\subsection{Architecture and ablation axes}

Every result reported here uses the same core configuration: $G=8$ groups,
alphabet $K_g=16$, $U=128$ units, $L=3$ layers, $D=256$, unit width $8$, trained
for 120 epochs with AdamW at a learning rate of $2\times10^{-3}$ and batch
size 1024.\footnote{The epoch count is set so that the number of optimiser steps
matches that of the smaller-batch configuration it replaced, making the two
training budgets comparable.} Only three axes vary, and each is named in the
configuration:

\begin{itemize}[leftmargin=*, itemsep=1pt, topsep=2pt]
  \item \textbf{encoder}, a single dense layer over flattened pixels, or a
        small CNN, both mapping into the same $\phi$ dimensionality;
  \item \textbf{core}, the sparse-joint core or a dense bottleneck of
        identical output dimensionality, used as a control;
  \item \textbf{$\lambda$}, the weight of the redundancy penalty of
        Equation~\ref{eq:redundancy}.
\end{itemize}

This matters for comparability. In an earlier state of this
work the image-classification experiments and the medical-imaging experiments
ran under different core configurations ($L=2, U=32, D=64$ against
$L=3, U=128, D=256$), which makes any cross-domain comparison between them
meaningless. Fixing the architecture is what licenses the generality claim in
Section~\ref{sec:results}.

\subsection{Validation}

We use stratified $k$-fold cross-validation with $k=5$ rather than the official
benchmark splits, pooling all official partitions first. Three consequences
follow, and all three are deliberate:

\begin{enumerate}[leftmargin=*, itemsep=1pt, topsep=2pt]
  \item every sample is evaluated exactly once, which matters on the smaller
        medical datasets where a single test split is noisy;
  \item the fold partitions are generated from a fixed seed and are
        \emph{identical across configurations}, so a comparison between two
        configurations is paired fold by fold;
  \item standardisation statistics are computed on the training fold only.
\end{enumerate}

Datasets are stratified-subsampled to a common ceiling of 17{,}000 samples, so
that MNIST does not receive four times the training budget of BloodMNIST. Where
a comparison is drawn between configurations, we report the paired difference
per fold, Cohen's $d$, and a Wilcoxon signed-rank $p$-value. With five folds
that test has little power; the effect size is the primary reading and the
$p$-value is reported for completeness rather than as a decision rule.

\begin{figure}[!t]
\centering
\resizebox{0.99\linewidth}{!}{%
\begin{tikzpicture}[font=\small, node distance=3mm]
\tikzset{
  box/.style={draw, rounded corners, align=center, minimum height=8mm, inner sep=3pt},
  ax/.style={draw, dashed, rounded corners, align=center, inner sep=3pt, fill=black!3},
}
\node[box] (data) {dataset\\(splits pooled)};
\node[box, right=7mm of data] (cap) {stratified\\sub-sample};
\node[box, right=7mm of cap] (fold) {stratified\\5-fold};
\node[ax, below=9mm of fold] (axes) {ablation axes\\encoder $\cdot$ core $\cdot$ $\lambda$};
\node[box, right=7mm of fold] (train) {train on $k-1$\\evaluate on $1$};
\node[box, right=7mm of train] (met) {metrics:\\task $+$ structure};

\draw[-{Latex[length=2mm]}] (data) -- (cap);
\draw[-{Latex[length=2mm]}] (cap) -- (fold);
\draw[-{Latex[length=2mm]}] (fold) -- (train);
\draw[-{Latex[length=2mm]}] (train) -- (met);
\draw[-{Latex[length=2mm]}] (axes) -- (train);
\node[below=1mm of axes, font=\scriptsize, align=center]
  {same partitions for every axis value\\$\Rightarrow$ fold-paired comparison};
\end{tikzpicture}%
}
\caption{The protocol. Fold partitions are fixed by seed and reused across every
configuration of a dataset, which is what makes the comparisons in
Table~\ref{tab:stats} paired rather than independent.}
\label{fig:protocol}
\end{figure}

\subsection{What we measure}

Task performance alone cannot distinguish a network whose groups carry
complementary information from one whose groups are redundant copies, since
both can classify equally well. We therefore report two families.

\textbf{Task.} Accuracy and macro one-vs-rest AUC, which are the metrics the
MedMNIST benchmark~\cite{yang2023medmnist} reports and therefore allow comparison against published
baselines; plus balanced accuracy, macro F1, Cohen's $\kappa$ and MCC. On
imbalanced medical data the gap between accuracy and balanced accuracy is itself
diagnostic: a single class accounts for roughly $67\%$ of DermaMNIST, and a
model that collapses onto it scores well on the former and poorly on the latter.

\textbf{Structure.} Let $z_g$ be the code of group $g$ and $y$ the label.
\begin{itemize}[leftmargin=*, itemsep=1pt, topsep=2pt]
  \item $\mathrm{NMI}(z_g; y)$, averaged over groups, which is how much any
        single group knows on its own;
  \item $\mathrm{NMI}(z_i; z_j)$, averaged over pairs, the \emph{redundancy};
        near zero means independent codes;
  \item \textbf{pair gain}, the quantity the architecture is built to produce:
        \begin{equation}
          \begin{aligned}
            \frac{1}{|E|}\sum_{(i,j)\in E} \Bigl[
              &\mathrm{NMI}\bigl((z_i,z_j); y\bigr)\\
              &- \max_{g \in \{i,j\}} \mathrm{NMI}(z_g;y) \Bigr],
          \end{aligned}
          \label{eq:pairgain}
        \end{equation}
        positive only when a connected pair carries more about the label
        \emph{together} than either member carries alone. This is the
        operational definition of compositionality we commit to;
  \item \textbf{VSA recovery}, the fraction of codes recovered by the
        bundle $\rightarrow$ unbind $\rightarrow$ cleanup round trip of
        Equations~\ref{eq:bundle} and~\ref{eq:cleanup}. It measures whether the
        program remains \emph{readable}, which any downstream composition
        between networks requires;
  \item alphabet utilisation and code perplexity, which detect a collapsed
        bottleneck that classifies well by ignoring most of its vocabulary.
\end{itemize}

\subsection{Reproducibility}

Each run writes its per-fold metrics, the aggregate, the dataset description and
a provenance record (commit hash, working-tree cleanliness, PyTorch version, GPU
model, timestamp) into a single JSON file. The tables and figures in this paper
are generated from those files by \texttt{report.py}, so no number here is
transcribed by hand. The code, the protocol and the result files are in the
project repository.
 depois da seção de trabalhos relacionados.
%
% Requer no preâmbulo (já presentes no template): graphicx, booktabs,
% tikz (arrows.meta, positioning), amsmath.
% =====================================================================

\section{The Sparse-Joint Core}
\label{sec:architecture}

A sparse-joint network replaces the dense bottleneck of a conventional encoder
with a \emph{graded program}: a representation that is simultaneously discrete,
as a tuple of per-group symbols, and continuous, as a hypervector that the
network can differentiate through. This section states the construction
precisely; Section~\ref{sec:protocol} describes how we test it.

\subsection{Groups and the scaffold}

Let $U$ be the number of hidden units per layer, $L$ the number of layers and
$G$ the number of \emph{groups}. A group is a sparse population of units
distributed across layers, fixed before training by a binary mask
$M \in \{0,1\}^{L \times G \times U}$, where $M_{l,g,u}=1$ means that group $g$
owns unit $u$ in layer $l$. We call $M$ the \emph{scaffold}.

Two groups are adjacent when they physically share a unit:
\begin{equation}
  (i,j) \in E \iff \exists\, l,u : M_{l,i,u} = M_{l,j,u} = 1 .
  \label{eq:edges}
\end{equation}
The edge set $E$ is therefore not a hyperparameter chosen alongside the
architecture; it is a \emph{consequence} of the scaffold, readable from the
weights themselves. This is what later allows $E$ to act as a compatibility
certificate between two independently trained networks: two networks with
isomorphic scaffolds expose the same relational skeleton, even though neither
has seen the other's data. Figure~\ref{fig:glue}(a) shows the graph used
throughout this work, for $G=8$.

Unit-to-unit transitions are routed through the scaffold: a connection between
units $u$ and $u'$ exists only if some group contains both. A group therefore
cannot read an arbitrary part of the hidden state, which is what makes the
sparsity structural rather than merely a regulariser.

\subsection{From activations to a program}

Write $h^{(l)} \in \mathbb{R}^{U \times W}$ for the hidden state at layer $l$,
each unit carrying a width-$W$ vector, with
$h^{(0)} = \mathrm{GELU}(W_{\text{in}}\, \phi(x))$ for an encoder $\phi$. Each
group pools the units it owns, across all layers:
\begin{equation}
  f_g = \frac{1}{L}\sum_{l=1}^{L}
        \frac{\sum_{u} M_{l,g,u}\, h^{(l)}_u}{\sum_{u} M_{l,g,u}},
  \qquad
  v_g = \frac{W_{\text{val}} f_g}{\lVert W_{\text{val}} f_g \rVert}.
  \label{eq:pool}
\end{equation}

Each group owns a fixed bipolar \emph{role} vector $r_g \in \{-1,+1\}^{D}$,
drawn once at initialisation. Binding is the Hadamard product, the standard
multiply-add vector-symbolic algebra: $r_g \odot v_g$ binds the group's value to
its identity, and because $r_g \odot r_g = \mathbf{1}$, binding again by the
same role unbinds it. Unary messages are written back to every unit the group
owns, and pairwise messages $r_i v_i \odot r_j v_j$ only to the units that $i$
and $j$ physically share, so interaction is confined to $E$ by construction.
The stack is then re-pooled, giving refined values $v'_g$.

The refined values are superposed into a single hypervector, and each group is
recovered from the superposition by unbinding with its own role:
\begin{equation}
  b = \frac{1}{\sqrt{G}}\sum_{g=1}^{G} r_g \odot v'_g,
  \qquad
  \tilde{v}_g = \sqrt{G}\, \bigl( b \odot r_g \bigr).
  \label{eq:bundle}
\end{equation}
Recovery is approximate: the other $G-1$ terms survive as noise whose magnitude
falls with $D$. We measure this directly (Section~\ref{sec:protocol},
\emph{VSA recovery}) rather than assuming it, because every downstream claim
about reading codes back out depends on it.

\subsection{Discretisation}

Each group owns a codebook $C_g \in \mathbb{R}^{K_g \times D}$ with its own
alphabet size $K_g$, so groups need not share a vocabulary. Cleanup is a cosine
lookup against that codebook,
\begin{equation}
  \ell_{g,c} = \frac{1}{\tau}
  \left\langle \frac{\tilde{v}_g}{\lVert \tilde{v}_g \rVert},
               \frac{C_{g,c}}{\lVert C_{g,c} \rVert} \right\rangle
  + m_{g,c},
  \label{eq:cleanup}
\end{equation}
where $m_{g,c} = -\infty$ for $c \geq K_g$ masks out fillers outside the group's
alphabet. The discrete code is $z_g = \arg\max_c \ell_{g,c}$, and gradients pass
through a straight-through estimator~\cite{bengio2013estimating,oord2017neural}, so the forward pass is exactly discrete
while the backward pass stays differentiable.

The exported \emph{program} is graded, concatenating a unary and a pairwise
bundle:
\begin{equation}
  \begin{aligned}
    P(x) &= \bigl[\, p_{\text{unary}} \;;\; p_{\text{pair}} \,\bigr],\\
    p_{\text{unary}} &= \tfrac{1}{\sqrt{G}}\textstyle\sum\nolimits_g r_g \odot \hat{v}_g,\\
    p_{\text{pair}} &= \tfrac{1}{\sqrt{|E|}}\textstyle\sum\nolimits_{(i,j) \in E}
       \hat{h}_i \odot \hat{h}_j,
  \end{aligned}
  \label{eq:program}
\end{equation}
where $\hat{h}_g = r_g \odot \hat{v}_g$ is the bound contribution of group $g$,
and $\hat{v}_g$ the quantised value of group $g$. The first term says what each
part is; the second says which parts co-occur, and only along scaffold edges.
Figure~\ref{fig:method} traces the whole path.

\subsection{Redundancy between groups}
\label{sec:penalty}

Nothing in Equations~\ref{eq:pool} to~\ref{eq:program} forces groups to encode
\emph{different} things. Trained only to classify, every group independently
converges on the label: the codes become near-copies of one another, and the
pairwise term of Equation~\ref{eq:program} carries nothing the unary term does
not already carry. We confirm this empirically in Section~\ref{sec:results}.

We therefore add a penalty on the mutual information between the soft code
distributions of every pair of groups. Writing $p_g \in \Delta^{K_g}$ for the
posterior of group $g$ on a batch, and $\hat{p}_{ij}$ for the empirical joint,
\begin{equation}
  \mathcal{L}_{\text{red}} =
  \frac{1}{G(G-1)} \sum_{i \neq j}
  \sum_{c,c'} \hat{p}_{ij}(c,c')
  \log \frac{\hat{p}_{ij}(c,c')}{\hat{p}_i(c)\,\hat{p}_j(c')},
  \label{eq:redundancy}
\end{equation}
added to the task loss with weight $\lambda$. Minimising it pushes the groups
towards statistically independent codes. The full objective is
$\mathcal{L} = \mathcal{L}_{\text{task}} + \mathcal{L}_{\text{VQ}}
+ \mathcal{L}_{\text{support}} + \lambda \mathcal{L}_{\text{red}}$.

% =====================================================================

\section{Experimental Protocol}
\label{sec:protocol}

The claims this architecture makes are about generality and about structure,
not only about accuracy, so the protocol has to be comparable across domains
and has to measure structure explicitly.

\subsection{Architecture and ablation axes}

Every result reported here uses the same core configuration: $G=8$ groups,
alphabet $K_g=16$, $U=128$ units, $L=3$ layers, $D=256$, unit width $8$, trained
for 120 epochs with AdamW at a learning rate of $2\times10^{-3}$ and batch
size 1024.\footnote{The epoch count is set so that the number of optimiser steps
matches that of the smaller-batch configuration it replaced, making the two
training budgets comparable.} Only three axes vary, and each is named in the
configuration:

\begin{itemize}[leftmargin=*, itemsep=1pt, topsep=2pt]
  \item \textbf{encoder}, a single dense layer over flattened pixels, or a
        small CNN, both mapping into the same $\phi$ dimensionality;
  \item \textbf{core}, the sparse-joint core or a dense bottleneck of
        identical output dimensionality, used as a control;
  \item \textbf{$\lambda$}, the weight of the redundancy penalty of
        Equation~\ref{eq:redundancy}.
\end{itemize}

This matters for comparability. In an earlier state of this
work the image-classification experiments and the medical-imaging experiments
ran under different core configurations ($L=2, U=32, D=64$ against
$L=3, U=128, D=256$), which makes any cross-domain comparison between them
meaningless. Fixing the architecture is what licenses the generality claim in
Section~\ref{sec:results}.

\subsection{Validation}

We use stratified $k$-fold cross-validation with $k=5$ rather than the official
benchmark splits, pooling all official partitions first. Three consequences
follow, and all three are deliberate:

\begin{enumerate}[leftmargin=*, itemsep=1pt, topsep=2pt]
  \item every sample is evaluated exactly once, which matters on the smaller
        medical datasets where a single test split is noisy;
  \item the fold partitions are generated from a fixed seed and are
        \emph{identical across configurations}, so a comparison between two
        configurations is paired fold by fold;
  \item standardisation statistics are computed on the training fold only.
\end{enumerate}

Datasets are stratified-subsampled to a common ceiling of 17{,}000 samples, so
that MNIST does not receive four times the training budget of BloodMNIST. Where
a comparison is drawn between configurations, we report the paired difference
per fold, Cohen's $d$, and a Wilcoxon signed-rank $p$-value. With five folds
that test has little power; the effect size is the primary reading and the
$p$-value is reported for completeness rather than as a decision rule.

\begin{figure}[!t]
\centering
\resizebox{0.99\linewidth}{!}{%
\begin{tikzpicture}[font=\small, node distance=3mm]
\tikzset{
  box/.style={draw, rounded corners, align=center, minimum height=8mm, inner sep=3pt},
  ax/.style={draw, dashed, rounded corners, align=center, inner sep=3pt, fill=black!3},
}
\node[box] (data) {dataset\\(splits pooled)};
\node[box, right=7mm of data] (cap) {stratified\\sub-sample};
\node[box, right=7mm of cap] (fold) {stratified\\5-fold};
\node[ax, below=9mm of fold] (axes) {ablation axes\\encoder $\cdot$ core $\cdot$ $\lambda$};
\node[box, right=7mm of fold] (train) {train on $k-1$\\evaluate on $1$};
\node[box, right=7mm of train] (met) {metrics:\\task $+$ structure};

\draw[-{Latex[length=2mm]}] (data) -- (cap);
\draw[-{Latex[length=2mm]}] (cap) -- (fold);
\draw[-{Latex[length=2mm]}] (fold) -- (train);
\draw[-{Latex[length=2mm]}] (train) -- (met);
\draw[-{Latex[length=2mm]}] (axes) -- (train);
\node[below=1mm of axes, font=\scriptsize, align=center]
  {same partitions for every axis value\\$\Rightarrow$ fold-paired comparison};
\end{tikzpicture}%
}
\caption{The protocol. Fold partitions are fixed by seed and reused across every
configuration of a dataset, which is what makes the comparisons in
Table~\ref{tab:stats} paired rather than independent.}
\label{fig:protocol}
\end{figure}

\subsection{What we measure}

Task performance alone cannot distinguish a network whose groups carry
complementary information from one whose groups are redundant copies, since
both can classify equally well. We therefore report two families.

\textbf{Task.} Accuracy and macro one-vs-rest AUC, which are the metrics the
MedMNIST benchmark~\cite{yang2023medmnist} reports and therefore allow comparison against published
baselines; plus balanced accuracy, macro F1, Cohen's $\kappa$ and MCC. On
imbalanced medical data the gap between accuracy and balanced accuracy is itself
diagnostic: a single class accounts for roughly $67\%$ of DermaMNIST, and a
model that collapses onto it scores well on the former and poorly on the latter.

\textbf{Structure.} Let $z_g$ be the code of group $g$ and $y$ the label.
\begin{itemize}[leftmargin=*, itemsep=1pt, topsep=2pt]
  \item $\mathrm{NMI}(z_g; y)$, averaged over groups, which is how much any
        single group knows on its own;
  \item $\mathrm{NMI}(z_i; z_j)$, averaged over pairs, the \emph{redundancy};
        near zero means independent codes;
  \item \textbf{pair gain}, the quantity the architecture is built to produce:
        \begin{equation}
          \begin{aligned}
            \frac{1}{|E|}\sum_{(i,j)\in E} \Bigl[
              &\mathrm{NMI}\bigl((z_i,z_j); y\bigr)\\
              &- \max_{g \in \{i,j\}} \mathrm{NMI}(z_g;y) \Bigr],
          \end{aligned}
          \label{eq:pairgain}
        \end{equation}
        positive only when a connected pair carries more about the label
        \emph{together} than either member carries alone. This is the
        operational definition of compositionality we commit to;
  \item \textbf{VSA recovery}, the fraction of codes recovered by the
        bundle $\rightarrow$ unbind $\rightarrow$ cleanup round trip of
        Equations~\ref{eq:bundle} and~\ref{eq:cleanup}. It measures whether the
        program remains \emph{readable}, which any downstream composition
        between networks requires;
  \item alphabet utilisation and code perplexity, which detect a collapsed
        bottleneck that classifies well by ignoring most of its vocabulary.
\end{itemize}

\subsection{Reproducibility}

Each run writes its per-fold metrics, the aggregate, the dataset description and
a provenance record (commit hash, working-tree cleanliness, PyTorch version, GPU
model, timestamp) into a single JSON file. The tables and figures in this paper
are generated from those files by \texttt{report.py}, so no number here is
transcribed by hand. The code, the protocol and the result files are in the
project repository.
 depois da seção de trabalhos relacionados.
%
% Requer no preâmbulo (já presentes no template): graphicx, booktabs,
% tikz (arrows.meta, positioning), amsmath.
% =====================================================================

\section{The Sparse-Joint Core}
\label{sec:architecture}

A sparse-joint network replaces the dense bottleneck of a conventional encoder
with a \emph{graded program}: a representation that is simultaneously discrete,
as a tuple of per-group symbols, and continuous, as a hypervector that the
network can differentiate through. This section states the construction
precisely; Section~\ref{sec:protocol} describes how we test it.

\subsection{Groups and the scaffold}

Let $U$ be the number of hidden units per layer, $L$ the number of layers and
$G$ the number of \emph{groups}. A group is a sparse population of units
distributed across layers, fixed before training by a binary mask
$M \in \{0,1\}^{L \times G \times U}$, where $M_{l,g,u}=1$ means that group $g$
owns unit $u$ in layer $l$. We call $M$ the \emph{scaffold}.

Two groups are adjacent when they physically share a unit:
\begin{equation}
  (i,j) \in E \iff \exists\, l,u : M_{l,i,u} = M_{l,j,u} = 1 .
  \label{eq:edges}
\end{equation}
The edge set $E$ is therefore not a hyperparameter chosen alongside the
architecture; it is a \emph{consequence} of the scaffold, readable from the
weights themselves. This is what later allows $E$ to act as a compatibility
certificate between two independently trained networks: two networks with
isomorphic scaffolds expose the same relational skeleton, even though neither
has seen the other's data. Figure~\ref{fig:glue}(a) shows the graph used
throughout this work, for $G=8$.

Unit-to-unit transitions are routed through the scaffold: a connection between
units $u$ and $u'$ exists only if some group contains both. A group therefore
cannot read an arbitrary part of the hidden state, which is what makes the
sparsity structural rather than merely a regulariser.

\subsection{From activations to a program}

Write $h^{(l)} \in \mathbb{R}^{U \times W}$ for the hidden state at layer $l$,
each unit carrying a width-$W$ vector, with
$h^{(0)} = \mathrm{GELU}(W_{\text{in}}\, \phi(x))$ for an encoder $\phi$. Each
group pools the units it owns, across all layers:
\begin{equation}
  f_g = \frac{1}{L}\sum_{l=1}^{L}
        \frac{\sum_{u} M_{l,g,u}\, h^{(l)}_u}{\sum_{u} M_{l,g,u}},
  \qquad
  v_g = \frac{W_{\text{val}} f_g}{\lVert W_{\text{val}} f_g \rVert}.
  \label{eq:pool}
\end{equation}

Each group owns a fixed bipolar \emph{role} vector $r_g \in \{-1,+1\}^{D}$,
drawn once at initialisation. Binding is the Hadamard product, the standard
multiply-add vector-symbolic algebra: $r_g \odot v_g$ binds the group's value to
its identity, and because $r_g \odot r_g = \mathbf{1}$, binding again by the
same role unbinds it. Unary messages are written back to every unit the group
owns, and pairwise messages $r_i v_i \odot r_j v_j$ only to the units that $i$
and $j$ physically share, so interaction is confined to $E$ by construction.
The stack is then re-pooled, giving refined values $v'_g$.

The refined values are superposed into a single hypervector, and each group is
recovered from the superposition by unbinding with its own role:
\begin{equation}
  b = \frac{1}{\sqrt{G}}\sum_{g=1}^{G} r_g \odot v'_g,
  \qquad
  \tilde{v}_g = \sqrt{G}\, \bigl( b \odot r_g \bigr).
  \label{eq:bundle}
\end{equation}
Recovery is approximate: the other $G-1$ terms survive as noise whose magnitude
falls with $D$. We measure this directly (Section~\ref{sec:protocol},
\emph{VSA recovery}) rather than assuming it, because every downstream claim
about reading codes back out depends on it.

\subsection{Discretisation}

Each group owns a codebook $C_g \in \mathbb{R}^{K_g \times D}$ with its own
alphabet size $K_g$, so groups need not share a vocabulary. Cleanup is a cosine
lookup against that codebook,
\begin{equation}
  \ell_{g,c} = \frac{1}{\tau}
  \left\langle \frac{\tilde{v}_g}{\lVert \tilde{v}_g \rVert},
               \frac{C_{g,c}}{\lVert C_{g,c} \rVert} \right\rangle
  + m_{g,c},
  \label{eq:cleanup}
\end{equation}
where $m_{g,c} = -\infty$ for $c \geq K_g$ masks out fillers outside the group's
alphabet. The discrete code is $z_g = \arg\max_c \ell_{g,c}$, and gradients pass
through a straight-through estimator~\cite{bengio2013estimating,oord2017neural}, so the forward pass is exactly discrete
while the backward pass stays differentiable.

The exported \emph{program} is graded, concatenating a unary and a pairwise
bundle:
\begin{equation}
  \begin{aligned}
    P(x) &= \bigl[\, p_{\text{unary}} \;;\; p_{\text{pair}} \,\bigr],\\
    p_{\text{unary}} &= \tfrac{1}{\sqrt{G}}\textstyle\sum\nolimits_g r_g \odot \hat{v}_g,\\
    p_{\text{pair}} &= \tfrac{1}{\sqrt{|E|}}\textstyle\sum\nolimits_{(i,j) \in E}
       \hat{h}_i \odot \hat{h}_j,
  \end{aligned}
  \label{eq:program}
\end{equation}
where $\hat{h}_g = r_g \odot \hat{v}_g$ is the bound contribution of group $g$,
and $\hat{v}_g$ the quantised value of group $g$. The first term says what each
part is; the second says which parts co-occur, and only along scaffold edges.
Figure~\ref{fig:method} traces the whole path.

\subsection{Redundancy between groups}
\label{sec:penalty}

Nothing in Equations~\ref{eq:pool} to~\ref{eq:program} forces groups to encode
\emph{different} things. Trained only to classify, every group independently
converges on the label: the codes become near-copies of one another, and the
pairwise term of Equation~\ref{eq:program} carries nothing the unary term does
not already carry. We confirm this empirically in Section~\ref{sec:results}.

We therefore add a penalty on the mutual information between the soft code
distributions of every pair of groups. Writing $p_g \in \Delta^{K_g}$ for the
posterior of group $g$ on a batch, and $\hat{p}_{ij}$ for the empirical joint,
\begin{equation}
  \mathcal{L}_{\text{red}} =
  \frac{1}{G(G-1)} \sum_{i \neq j}
  \sum_{c,c'} \hat{p}_{ij}(c,c')
  \log \frac{\hat{p}_{ij}(c,c')}{\hat{p}_i(c)\,\hat{p}_j(c')},
  \label{eq:redundancy}
\end{equation}
added to the task loss with weight $\lambda$. Minimising it pushes the groups
towards statistically independent codes. The full objective is
$\mathcal{L} = \mathcal{L}_{\text{task}} + \mathcal{L}_{\text{VQ}}
+ \mathcal{L}_{\text{support}} + \lambda \mathcal{L}_{\text{red}}$.

% =====================================================================

\section{Experimental Protocol}
\label{sec:protocol}

The claims this architecture makes are about generality and about structure,
not only about accuracy, so the protocol has to be comparable across domains
and has to measure structure explicitly.

\subsection{Architecture and ablation axes}

Every result reported here uses the same core configuration: $G=8$ groups,
alphabet $K_g=16$, $U=128$ units, $L=3$ layers, $D=256$, unit width $8$, trained
for 120 epochs with AdamW at a learning rate of $2\times10^{-3}$ and batch
size 1024.\footnote{The epoch count is set so that the number of optimiser steps
matches that of the smaller-batch configuration it replaced, making the two
training budgets comparable.} Only three axes vary, and each is named in the
configuration:

\begin{itemize}[leftmargin=*, itemsep=1pt, topsep=2pt]
  \item \textbf{encoder}, a single dense layer over flattened pixels, or a
        small CNN, both mapping into the same $\phi$ dimensionality;
  \item \textbf{core}, the sparse-joint core or a dense bottleneck of
        identical output dimensionality, used as a control;
  \item \textbf{$\lambda$}, the weight of the redundancy penalty of
        Equation~\ref{eq:redundancy}.
\end{itemize}

This matters for comparability. In an earlier state of this
work the image-classification experiments and the medical-imaging experiments
ran under different core configurations ($L=2, U=32, D=64$ against
$L=3, U=128, D=256$), which makes any cross-domain comparison between them
meaningless. Fixing the architecture is what licenses the generality claim in
Section~\ref{sec:results}.

\subsection{Validation}

We use stratified $k$-fold cross-validation with $k=5$ rather than the official
benchmark splits, pooling all official partitions first. Three consequences
follow, and all three are deliberate:

\begin{enumerate}[leftmargin=*, itemsep=1pt, topsep=2pt]
  \item every sample is evaluated exactly once, which matters on the smaller
        medical datasets where a single test split is noisy;
  \item the fold partitions are generated from a fixed seed and are
        \emph{identical across configurations}, so a comparison between two
        configurations is paired fold by fold;
  \item standardisation statistics are computed on the training fold only.
\end{enumerate}

Datasets are stratified-subsampled to a common ceiling of 17{,}000 samples, so
that MNIST does not receive four times the training budget of BloodMNIST. Where
a comparison is drawn between configurations, we report the paired difference
per fold, Cohen's $d$, and a Wilcoxon signed-rank $p$-value. With five folds
that test has little power; the effect size is the primary reading and the
$p$-value is reported for completeness rather than as a decision rule.

\begin{figure}[!t]
\centering
\resizebox{0.99\linewidth}{!}{%
\begin{tikzpicture}[font=\small, node distance=3mm]
\tikzset{
  box/.style={draw, rounded corners, align=center, minimum height=8mm, inner sep=3pt},
  ax/.style={draw, dashed, rounded corners, align=center, inner sep=3pt, fill=black!3},
}
\node[box] (data) {dataset\\(splits pooled)};
\node[box, right=7mm of data] (cap) {stratified\\sub-sample};
\node[box, right=7mm of cap] (fold) {stratified\\5-fold};
\node[ax, below=9mm of fold] (axes) {ablation axes\\encoder $\cdot$ core $\cdot$ $\lambda$};
\node[box, right=7mm of fold] (train) {train on $k-1$\\evaluate on $1$};
\node[box, right=7mm of train] (met) {metrics:\\task $+$ structure};

\draw[-{Latex[length=2mm]}] (data) -- (cap);
\draw[-{Latex[length=2mm]}] (cap) -- (fold);
\draw[-{Latex[length=2mm]}] (fold) -- (train);
\draw[-{Latex[length=2mm]}] (train) -- (met);
\draw[-{Latex[length=2mm]}] (axes) -- (train);
\node[below=1mm of axes, font=\scriptsize, align=center]
  {same partitions for every axis value\\$\Rightarrow$ fold-paired comparison};
\end{tikzpicture}%
}
\caption{The protocol. Fold partitions are fixed by seed and reused across every
configuration of a dataset, which is what makes the comparisons in
Table~\ref{tab:stats} paired rather than independent.}
\label{fig:protocol}
\end{figure}

\subsection{What we measure}

Task performance alone cannot distinguish a network whose groups carry
complementary information from one whose groups are redundant copies, since
both can classify equally well. We therefore report two families.

\textbf{Task.} Accuracy and macro one-vs-rest AUC, which are the metrics the
MedMNIST benchmark~\cite{yang2023medmnist} reports and therefore allow comparison against published
baselines; plus balanced accuracy, macro F1, Cohen's $\kappa$ and MCC. On
imbalanced medical data the gap between accuracy and balanced accuracy is itself
diagnostic: a single class accounts for roughly $67\%$ of DermaMNIST, and a
model that collapses onto it scores well on the former and poorly on the latter.

\textbf{Structure.} Let $z_g$ be the code of group $g$ and $y$ the label.
\begin{itemize}[leftmargin=*, itemsep=1pt, topsep=2pt]
  \item $\mathrm{NMI}(z_g; y)$, averaged over groups, which is how much any
        single group knows on its own;
  \item $\mathrm{NMI}(z_i; z_j)$, averaged over pairs, the \emph{redundancy};
        near zero means independent codes;
  \item \textbf{pair gain}, the quantity the architecture is built to produce:
        \begin{equation}
          \begin{aligned}
            \frac{1}{|E|}\sum_{(i,j)\in E} \Bigl[
              &\mathrm{NMI}\bigl((z_i,z_j); y\bigr)\\
              &- \max_{g \in \{i,j\}} \mathrm{NMI}(z_g;y) \Bigr],
          \end{aligned}
          \label{eq:pairgain}
        \end{equation}
        positive only when a connected pair carries more about the label
        \emph{together} than either member carries alone. This is the
        operational definition of compositionality we commit to;
  \item \textbf{VSA recovery}, the fraction of codes recovered by the
        bundle $\rightarrow$ unbind $\rightarrow$ cleanup round trip of
        Equations~\ref{eq:bundle} and~\ref{eq:cleanup}. It measures whether the
        program remains \emph{readable}, which any downstream composition
        between networks requires;
  \item alphabet utilisation and code perplexity, which detect a collapsed
        bottleneck that classifies well by ignoring most of its vocabulary.
\end{itemize}

\subsection{Reproducibility}

Each run writes its per-fold metrics, the aggregate, the dataset description and
a provenance record (commit hash, working-tree cleanliness, PyTorch version, GPU
model, timestamp) into a single JSON file. The tables and figures in this paper
are generated from those files by \texttt{report.py}, so no number here is
transcribed by hand. The code, the protocol and the result files are in the
project repository.
 depois da seção de trabalhos relacionados.
%
% Requer no preâmbulo (já presentes no template): graphicx, booktabs,
% tikz (arrows.meta, positioning), amsmath.
% =====================================================================

\section{The Sparse-Joint Core}
\label{sec:architecture}

A sparse-joint network replaces the dense bottleneck of a conventional encoder
with a \emph{graded program}: a representation that is simultaneously discrete,
as a tuple of per-group symbols, and continuous, as a hypervector that the
network can differentiate through. This section states the construction
precisely; Section~\ref{sec:protocol} describes how we test it.

\subsection{Groups and the scaffold}

Let $U$ be the number of hidden units per layer, $L$ the number of layers and
$G$ the number of \emph{groups}. A group is a sparse population of units
distributed across layers, fixed before training by a binary mask
$M \in \{0,1\}^{L \times G \times U}$, where $M_{l,g,u}=1$ means that group $g$
owns unit $u$ in layer $l$. We call $M$ the \emph{scaffold}.

Two groups are adjacent when they physically share a unit:
\begin{equation}
  (i,j) \in E \iff \exists\, l,u : M_{l,i,u} = M_{l,j,u} = 1 .
  \label{eq:edges}
\end{equation}
The edge set $E$ is therefore not a hyperparameter chosen alongside the
architecture; it is a \emph{consequence} of the scaffold, readable from the
weights themselves. This is what later allows $E$ to act as a compatibility
certificate between two independently trained networks: two networks with
isomorphic scaffolds expose the same relational skeleton, even though neither
has seen the other's data. Figure~\ref{fig:glue}(a) shows the graph used
throughout this work, for $G=8$.

Unit-to-unit transitions are routed through the scaffold: a connection between
units $u$ and $u'$ exists only if some group contains both. A group therefore
cannot read an arbitrary part of the hidden state, which is what makes the
sparsity structural rather than merely a regulariser.

\subsection{From activations to a program}

Write $h^{(l)} \in \mathbb{R}^{U \times W}$ for the hidden state at layer $l$,
each unit carrying a width-$W$ vector, with
$h^{(0)} = \mathrm{GELU}(W_{\text{in}}\, \phi(x))$ for an encoder $\phi$. Each
group pools the units it owns, across all layers:
\begin{equation}
  f_g = \frac{1}{L}\sum_{l=1}^{L}
        \frac{\sum_{u} M_{l,g,u}\, h^{(l)}_u}{\sum_{u} M_{l,g,u}},
  \qquad
  v_g = \frac{W_{\text{val}} f_g}{\lVert W_{\text{val}} f_g \rVert}.
  \label{eq:pool}
\end{equation}

Each group owns a fixed bipolar \emph{role} vector $r_g \in \{-1,+1\}^{D}$,
drawn once at initialisation. Binding is the Hadamard product, the standard
multiply-add vector-symbolic algebra: $r_g \odot v_g$ binds the group's value to
its identity, and because $r_g \odot r_g = \mathbf{1}$, binding again by the
same role unbinds it. Unary messages are written back to every unit the group
owns, and pairwise messages $r_i v_i \odot r_j v_j$ only to the units that $i$
and $j$ physically share, so interaction is confined to $E$ by construction.
The stack is then re-pooled, giving refined values $v'_g$.

The refined values are superposed into a single hypervector, and each group is
recovered from the superposition by unbinding with its own role:
\begin{equation}
  b = \frac{1}{\sqrt{G}}\sum_{g=1}^{G} r_g \odot v'_g,
  \qquad
  \tilde{v}_g = \sqrt{G}\, \bigl( b \odot r_g \bigr).
  \label{eq:bundle}
\end{equation}
Recovery is approximate: the other $G-1$ terms survive as noise whose magnitude
falls with $D$. We measure this directly (Section~\ref{sec:protocol},
\emph{VSA recovery}) rather than assuming it, because every downstream claim
about reading codes back out depends on it.

\subsection{Discretisation}

Each group owns a codebook $C_g \in \mathbb{R}^{K_g \times D}$ with its own
alphabet size $K_g$, so groups need not share a vocabulary. Cleanup is a cosine
lookup against that codebook,
\begin{equation}
  \ell_{g,c} = \frac{1}{\tau}
  \left\langle \frac{\tilde{v}_g}{\lVert \tilde{v}_g \rVert},
               \frac{C_{g,c}}{\lVert C_{g,c} \rVert} \right\rangle
  + m_{g,c},
  \label{eq:cleanup}
\end{equation}
where $m_{g,c} = -\infty$ for $c \geq K_g$ masks out fillers outside the group's
alphabet. The discrete code is $z_g = \arg\max_c \ell_{g,c}$, and gradients pass
through a straight-through estimator~\cite{bengio2013estimating,oord2017neural}, so the forward pass is exactly discrete
while the backward pass stays differentiable.

The exported \emph{program} is graded, concatenating a unary and a pairwise
bundle:
\begin{equation}
  \begin{aligned}
    P(x) &= \bigl[\, p_{\text{unary}} \;;\; p_{\text{pair}} \,\bigr],\\
    p_{\text{unary}} &= \tfrac{1}{\sqrt{G}}\textstyle\sum\nolimits_g r_g \odot \hat{v}_g,\\
    p_{\text{pair}} &= \tfrac{1}{\sqrt{|E|}}\textstyle\sum\nolimits_{(i,j) \in E}
       \hat{h}_i \odot \hat{h}_j,
  \end{aligned}
  \label{eq:program}
\end{equation}
where $\hat{h}_g = r_g \odot \hat{v}_g$ is the bound contribution of group $g$,
and $\hat{v}_g$ the quantised value of group $g$. The first term says what each
part is; the second says which parts co-occur, and only along scaffold edges.
Figure~\ref{fig:method} traces the whole path.

\subsection{Redundancy between groups}
\label{sec:penalty}

Nothing in Equations~\ref{eq:pool} to~\ref{eq:program} forces groups to encode
\emph{different} things. Trained only to classify, every group independently
converges on the label: the codes become near-copies of one another, and the
pairwise term of Equation~\ref{eq:program} carries nothing the unary term does
not already carry. We confirm this empirically in Section~\ref{sec:results}.

We therefore add a penalty on the mutual information between the soft code
distributions of every pair of groups. Writing $p_g \in \Delta^{K_g}$ for the
posterior of group $g$ on a batch, and $\hat{p}_{ij}$ for the empirical joint,
\begin{equation}
  \mathcal{L}_{\text{red}} =
  \frac{1}{G(G-1)} \sum_{i \neq j}
  \sum_{c,c'} \hat{p}_{ij}(c,c')
  \log \frac{\hat{p}_{ij}(c,c')}{\hat{p}_i(c)\,\hat{p}_j(c')},
  \label{eq:redundancy}
\end{equation}
added to the task loss with weight $\lambda$. Minimising it pushes the groups
towards statistically independent codes. The full objective is
$\mathcal{L} = \mathcal{L}_{\text{task}} + \mathcal{L}_{\text{VQ}}
+ \mathcal{L}_{\text{support}} + \lambda \mathcal{L}_{\text{red}}$.

% =====================================================================

\section{Experimental Protocol}
\label{sec:protocol}

The claims this architecture makes are about generality and about structure,
not only about accuracy, so the protocol has to be comparable across domains
and has to measure structure explicitly.

\subsection{Architecture and ablation axes}

Every result reported here uses the same core configuration: $G=8$ groups,
alphabet $K_g=16$, $U=128$ units, $L=3$ layers, $D=256$, unit width $8$, trained
for 120 epochs with AdamW at a learning rate of $2\times10^{-3}$ and batch
size 1024.\footnote{The epoch count is set so that the number of optimiser steps
matches that of the smaller-batch configuration it replaced, making the two
training budgets comparable.} Only three axes vary, and each is named in the
configuration:

\begin{itemize}[leftmargin=*, itemsep=1pt, topsep=2pt]
  \item \textbf{encoder}, a single dense layer over flattened pixels, or a
        small CNN, both mapping into the same $\phi$ dimensionality;
  \item \textbf{core}, the sparse-joint core or a dense bottleneck of
        identical output dimensionality, used as a control;
  \item \textbf{$\lambda$}, the weight of the redundancy penalty of
        Equation~\ref{eq:redundancy}.
\end{itemize}

This matters for comparability. In an earlier state of this
work the image-classification experiments and the medical-imaging experiments
ran under different core configurations ($L=2, U=32, D=64$ against
$L=3, U=128, D=256$), which makes any cross-domain comparison between them
meaningless. Fixing the architecture is what licenses the generality claim in
Section~\ref{sec:results}.

\subsection{Validation}

We use stratified $k$-fold cross-validation with $k=5$ rather than the official
benchmark splits, pooling all official partitions first. Three consequences
follow, and all three are deliberate:

\begin{enumerate}[leftmargin=*, itemsep=1pt, topsep=2pt]
  \item every sample is evaluated exactly once, which matters on the smaller
        medical datasets where a single test split is noisy;
  \item the fold partitions are generated from a fixed seed and are
        \emph{identical across configurations}, so a comparison between two
        configurations is paired fold by fold;
  \item standardisation statistics are computed on the training fold only.
\end{enumerate}

Datasets are stratified-subsampled to a common ceiling of 17{,}000 samples, so
that MNIST does not receive four times the training budget of BloodMNIST. Where
a comparison is drawn between configurations, we report the paired difference
per fold, Cohen's $d$, and a Wilcoxon signed-rank $p$-value. With five folds
that test has little power; the effect size is the primary reading and the
$p$-value is reported for completeness rather than as a decision rule.

\begin{figure}[!t]
\centering
\resizebox{0.99\linewidth}{!}{%
\begin{tikzpicture}[font=\small, node distance=3mm]
\tikzset{
  box/.style={draw, rounded corners, align=center, minimum height=8mm, inner sep=3pt},
  ax/.style={draw, dashed, rounded corners, align=center, inner sep=3pt, fill=black!3},
}
\node[box] (data) {dataset\\(splits pooled)};
\node[box, right=7mm of data] (cap) {stratified\\sub-sample};
\node[box, right=7mm of cap] (fold) {stratified\\5-fold};
\node[ax, below=9mm of fold] (axes) {ablation axes\\encoder $\cdot$ core $\cdot$ $\lambda$};
\node[box, right=7mm of fold] (train) {train on $k-1$\\evaluate on $1$};
\node[box, right=7mm of train] (met) {metrics:\\task $+$ structure};

\draw[-{Latex[length=2mm]}] (data) -- (cap);
\draw[-{Latex[length=2mm]}] (cap) -- (fold);
\draw[-{Latex[length=2mm]}] (fold) -- (train);
\draw[-{Latex[length=2mm]}] (train) -- (met);
\draw[-{Latex[length=2mm]}] (axes) -- (train);
\node[below=1mm of axes, font=\scriptsize, align=center]
  {same partitions for every axis value\\$\Rightarrow$ fold-paired comparison};
\end{tikzpicture}%
}
\caption{The protocol. Fold partitions are fixed by seed and reused across every
configuration of a dataset, which is what makes the comparisons in
Table~\ref{tab:stats} paired rather than independent.}
\label{fig:protocol}
\end{figure}

\subsection{What we measure}

Task performance alone cannot distinguish a network whose groups carry
complementary information from one whose groups are redundant copies, since
both can classify equally well. We therefore report two families.

\textbf{Task.} Accuracy and macro one-vs-rest AUC, which are the metrics the
MedMNIST benchmark~\cite{yang2023medmnist} reports and therefore allow comparison against published
baselines; plus balanced accuracy, macro F1, Cohen's $\kappa$ and MCC. On
imbalanced medical data the gap between accuracy and balanced accuracy is itself
diagnostic: a single class accounts for roughly $67\%$ of DermaMNIST, and a
model that collapses onto it scores well on the former and poorly on the latter.

\textbf{Structure.} Let $z_g$ be the code of group $g$ and $y$ the label.
\begin{itemize}[leftmargin=*, itemsep=1pt, topsep=2pt]
  \item $\mathrm{NMI}(z_g; y)$, averaged over groups, which is how much any
        single group knows on its own;
  \item $\mathrm{NMI}(z_i; z_j)$, averaged over pairs, the \emph{redundancy};
        near zero means independent codes;
  \item \textbf{pair gain}, the quantity the architecture is built to produce:
        \begin{equation}
          \begin{aligned}
            \frac{1}{|E|}\sum_{(i,j)\in E} \Bigl[
              &\mathrm{NMI}\bigl((z_i,z_j); y\bigr)\\
              &- \max_{g \in \{i,j\}} \mathrm{NMI}(z_g;y) \Bigr],
          \end{aligned}
          \label{eq:pairgain}
        \end{equation}
        positive only when a connected pair carries more about the label
        \emph{together} than either member carries alone. This is the
        operational definition of compositionality we commit to;
  \item \textbf{VSA recovery}, the fraction of codes recovered by the
        bundle $\rightarrow$ unbind $\rightarrow$ cleanup round trip of
        Equations~\ref{eq:bundle} and~\ref{eq:cleanup}. It measures whether the
        program remains \emph{readable}, which any downstream composition
        between networks requires;
  \item alphabet utilisation and code perplexity, which detect a collapsed
        bottleneck that classifies well by ignoring most of its vocabulary.
\end{itemize}

\subsection{Reproducibility}

Each run writes its per-fold metrics, the aggregate, the dataset description and
a provenance record (commit hash, working-tree cleanliness, PyTorch version, GPU
model, timestamp) into a single JSON file. The tables and figures in this paper
are generated from those files by \texttt{report.py}, so no number here is
transcribed by hand. The code, the protocol and the result files are in the
project repository.
